\documentclass[12pt,a4paper]{article}
\usepackage[margin=2.5cm]{geometry}
\usepackage{lmodern}
\usepackage[T1]{fontenc}
\usepackage{amsmath,amssymb}
\usepackage{booktabs,array,tabularx}
\usepackage{enumitem}
\usepackage{xcolor}
\usepackage[most]{tcolorbox}
\usepackage{titlesec}
\usepackage[colorlinks=true,linkcolor=black,urlcolor=blue!50!black]{hyperref}

\definecolor{headblue}{RGB}{38,70,120}
\titleformat{\section}{\Large\bfseries\color{headblue}}{\thesection}{0.8em}{}[{\color{black}\titlerule}]
\titleformat{\subsection}{\large\bfseries}{\thesubsection}{0.8em}{}
\setlength{\parindent}{0pt}
\setlength{\parskip}{6pt}

\newcommand{\mdesc}[1]{\par{\textit{Meaning:} #1}\par}
\newtcolorbox{objbox}{colback=blue!3,colframe=headblue!60,boxrule=0.5pt,arc=2mm,left=2mm,right=2mm,top=1mm,bottom=1mm}

\begin{document}

\begin{center}
{\LARGE\bfseries\color{headblue} EV Charging Network Optimisation}\\[6pt]
{\large Formulation 3: EV $\rightarrow$ charger (FCS / MCS) assignment + mobile charger routing\\ with a continuous income attribute (market segmentation)}\\[4pt]
{\normalsize CP 320 Operations Research for Mobility Management, course project}\\[2pt]
{\normalsize 21 November 2026}
\end{center}
\vspace{4pt}

%=====================================================================
\section{Problem Statement}

Given a small road network, the network operator wants to assign each EV with its respective demand either to a fixed charging station or a mobile charging station.

\hspace{1.5em}During the time interval $[0,T]$ a set of parked EVs raise charging requests. Each request must be served by exactly one resource: either the EV drives itself to a fixed charging station (FCS, 2 guns of 22\,kW each), or a mobile charging van (MCS) leaves the depot, drives to the EV, charges it (one EV at a time) and moves on, returning to the depot by $T$.

Customers are not put into classes beforehand. Every EV $i$ carries an \textbf{income index} $\theta_i\in[0,1]$: $\theta_i=0$ is the lowest income and $\theta_i=1$ the highest. Income changes how a customer feels a rupee and a minute:
\begin{itemize}[itemsep=1pt]
  \item a lower-income customer feels every rupee more (higher price sensitivity), so the cheaper FCS appeals to them;
  \item a higher-income customer values time more (higher value of time) and is less willing to queue, so a van that comes to the car appeals to them.
\end{itemize}
Who ends up at an FCS and who gets a van is not fixed in advance. It follows from these sensitivities, the operator's costs and the limited guns and vans.

We decide -
\begin{enumerate}[itemsep=2pt]
  \item Who serves each EV.
  \item The order in which each van visits its EVs.
  \item The time every charging session starts.
\end{enumerate}

\begin{objbox}
{\color{headblue}\textbf{Our Objective}}\\
Minimise the operator's cost plus the cost that customers perceive, where each customer weighs money and time according to their own income.
\end{objbox}

%=====================================================================
\section{Assumptions}
\begin{enumerate}[itemsep=2pt]
  \item All requests of the planning period are known at time 0; request $i$ becomes active at $r_i$ (request release time). All times are non-negative real numbers, measured in minutes, in $[0,T]$.
  \item An EV served by an FCS drives there itself (one way) on its own battery, waits for a free gun, and stays until charged. An EV can only use an FCS within its remaining range.
  \item Customer $i$ waits at most $W_i$ minutes at a station for a free gun; $W_i$ depends on income.
  \item An EV served by an MCS stays parked at its location; the van comes to it.
  \item Each FCS gun and each MCS charges one EV at a time, with no pre-emption.
  \item A van's driving energy and delivered energy come from the same battery; there is no recharging of the van during the peak.
  \item Every request must be fully charged by $T$ (hard).
  \item Charging power is constant.
  \item Every customer has an income index $\theta_i\in[0,1]$, known in advance. The $\theta_i$ are drawn independently from one distribution on $[0,1]$.
  \item Income acts only through three sensitivities, which are known functions of $\theta$: the money sensitivity $\mu(\theta)$ is non-increasing, the value of time $c^{a}(\theta)$ of each activity $a$ is non-decreasing and the queue tolerance $W(\theta)$ is non-increasing.
  \item Customers pay a tariff per kWh: $p^{\mathrm F}$ at an FCS and $p^{\mathrm M}\ge p^{\mathrm F}$ for the van service (a premium for charging at the door). An FCS user also pays for driving to the station.
  \item The operator minimises its own cost plus the customers' perceived cost. Tariff revenue is not counted as a gain for the operator.
\end{enumerate}

%=====================================================================
\newpage
\section{Sets, parameters and decision variables}

\subsection*{Sets}
\begin{tabularx}{\textwidth}{@{}>{\raggedright\arraybackslash}p{4.6cm}X@{}}
\toprule
\textbf{Notation} & \textbf{Description} \\
\midrule
$I=\{1,\dots,n\}$ & Set of EV charging requests. \\[2pt]
$F$,\ \ $|F|=12$ (or $24$) & Set of fixed charging stations (FCS). Dataset: 12 stations in the 100-EV files, 24 in the 200-EV files. \\[2pt]
$G_f$,\ \ $|G_f|=2$ & Set of guns (charging points) at FCS $f$, with $|G_f|=2$. Define $G=\bigcup_{f\in F}G_f$, and let $f(g)$ denote the station associated with gun $g$. Dataset: $|G|=24$ (or $48$). \\[2pt]
$F_i=\{f\in F: d_{if}\le R_i\}$ & Set of FCS that EV $i$ can reach given its remaining range. \\[2pt]
$G_i=\bigcup_{f\in F_i}G_f$ & Set of guns that EV $i$ can use. \\[2pt]
$M$,\ \ $|M|=5$ (or $12$) & Set of mobile charging vans. Dataset: 5 vans in the 100-EV files, 12 in the 200-EV files, for every instance size $n$. \\[2pt]
$0,\ n+1$ & Start and end depot nodes, respectively, representing the same physical location. (assumed single depot) \\[2pt]
$N=\{0\}\cup I\cup\{n+1\}$ & Set of routing nodes. \\[2pt]
$A$ & Set of allowed van-routing arcs, defined as \newline
$A=\{(i,j): i\in\{0\}\cup I,$\newline
\hspace*{3.2em}$j\in I\cup\{n+1\},$\newline
\hspace*{3.2em}$i\ne j,\ (i,j)\ne(0,n+1)\}$. \\[2pt]
$P_g=\{(i,j)\in I^2: i<j,\ g\in G_i\cap G_j\}$ & Set of EV pairs that could be assigned to the same gun $g$ (and must then be sequenced). \\[2pt]
$G^{+}=\{g\in G: g\in G_i \text{ for some } i\in I\}$ & Set of guns that at least one EV can reach. \\[2pt]
$\mathcal A=\{\mathrm D,\mathrm Q,\mathrm C,\mathrm V\}$ & Customer activities that take time: driving to an FCS, queueing at an FCS, charging at an FCS (driver present for all three), and the car parked while it waits for and is charged by a van (driver free). \\
\bottomrule
\end{tabularx}

\newpage
\subsection*{Parameters}
{\small
\begin{tabularx}{\textwidth}{@{}>{\raggedright\arraybackslash}p{4.6cm}X@{}}
\toprule
\textbf{Notation} & \textbf{Description} \\
\midrule
$E_i\in[10,40]$ & Energy demand of EV $i$ (kWh). Dataset: Solomon DEMAND (1--50) mapped linearly, $E_i=10+30\,(\text{DEMAND}_i-1)/49$. \\[2pt]
$r_i$ & Request (release) time (minutes). [$r_i=0$ for static demand] Dataset: Solomon READY TIME, unchanged. \\[2pt]
$R_i$ & Remaining driving range of EV $i$ (km). Dataset: uniform on $[5,40]$; for 15\,\% of the EVs (``stranded'') $R_i$ is 20--80\,\% of the distance to their nearest FCS, so $F_i=\emptyset$ and only a van can serve them. \\[2pt]
$P^{\mathrm F}=22,\ P^{\mathrm M}=30$ & Charging power of an FCS gun and an MCS (kW). \\[2pt]
$\tau^{\mathrm F}_i=\dfrac{60\,E_i}{P^{\mathrm F}},\ \tau^{\mathrm M}_i=\dfrac{60\,E_i}{P^{\mathrm M}}$ & Charging duration of EV $i$ at an FCS and by an MCS, respectively (minutes). Dataset: $\tau^{\mathrm F}_i=60E_i/22\in[27,109]$, $\tau^{\mathrm M}_i=60E_i/30\in[20,80]$. \\[2pt]
$d_{if},\ t_{if}=\dfrac{60\,d_{if}}{v}$ & Shortest-path distance (km) and travel time (minutes) from EV $i$ to FCS $f$. \\[2pt]
$d_{ij},\ t_{ij}=\dfrac{60\,d_{ij}}{v}$ & Shortest-path distance (km) and travel time (minutes) between routing nodes. Dataset: $d=1.3\times0.2\times$ (straight-line distance in Solomon grid units), i.e.\ 1 grid unit $=0.2$\,km and a road detour factor of 1.3. \\[2pt]
$v=15.6$ & Travel speed (km/h), assumed identical for all vehicles. Dataset: $v=60\times0.2\times1.3=15.6$\,km/h, so that $t$ in minutes equals the straight-line grid distance (Solomon's rule). \\[2pt]
$B_m=150$ & Usable battery capacity of MCS $m$ (kWh). \\[2pt]
$Q_m=6$ & Maximum number of EVs MCS $m$ may serve during the planning period. \\[2pt]
$\beta=0.25$ & Van driving energy consumption (kWh/km). \\[2pt]
$\eta=0.9$ & MCS charging efficiency (battery to EV), $\eta\in(0,1]$. \\[2pt]
$T$ & Length of the planning period (minutes). Dataset: Solomon's depot due date, unchanged. 100-EV files: C1 1236, C2 3390, R1 230, R2 1000, RC1 240, RC2 960. 200-EV files: C1 1351, C2 3598, R1 634, R2 2535, RC1 634, RC2 2535. \\[2pt]
$c^{\mathrm{EV}}=4,\ c^{\mathrm M}=15$ & Cost per kilometre of EV driving and van driving, respectively (Rs/km). \\[2pt]
$c^{\mathrm F}_m=1500$ & Fixed dispatch cost of van $m$ (driver, depreciation) (Rs). \\[2pt]
$\pi^{\mathrm F}=8.7,\ \pi^{\mathrm M}=10.7$ & Operator's cost per kWh of grid energy at an FCS and of energy stored in a van battery, respectively (Rs/kWh). \\[2pt]
$M^{\mathrm M}_{ij}=T+\tau^{\mathrm M}_i+t_{ij}$ & Big-$M$ constant for van sequencing constraints (minutes). \\[2pt]
$M^{\mathrm F}_i=T+\tau^{\mathrm F}_i$ & Big-$M$ constant for gun sequencing constraints (minutes). \\
\bottomrule
\end{tabularx}}

\newpage
\subsection*{Income parameters (new in formulation 3)}
{\small
\begin{tabularx}{\textwidth}{@{}>{\raggedright\arraybackslash}p{4.6cm}X@{}}
\toprule
\textbf{Notation} & \textbf{Description} \\
\midrule
$\theta_i\in[0,1]$ & Income index of the owner of EV $i$: 0 = lowest income, 1 = highest. Dataset: drawn independently from a Beta$(a_\theta,b_\theta)$ distribution; the same EV number has the same $\theta_i$ in every file of a Solomon class (positions, demands and ranges are shared within a class too). \\[2pt]
$\mu_i=\mu(\theta_i)>0$ & Money sensitivity of customer $i$: perceived cost per rupee paid (Rs/Rs). $\mu(\cdot)$ is non-increasing: a poorer customer feels a rupee more. \\[2pt]
$c^{a}_i=c^{a}(\theta_i)\ge0$,\ \ $a\in\mathcal A$ & Value of time of customer $i$ while driving to an FCS ($a=\mathrm D$), queueing at an FCS ($\mathrm Q$), charging at an FCS ($\mathrm C$), and while the car waits for and is charged by a van ($\mathrm V$) (Rs/minute). Each $c^{a}(\cdot)$ is non-decreasing: a richer customer values time more. At every income, queueing is worse than driving, $c^{\mathrm Q}(\theta)>c^{\mathrm D}(\theta)$, and a minute with the driver free costs less than a minute at the station, $c^{\mathrm V}(\theta)\le c^{\mathrm C}(\theta)$, with the gap not shrinking as income rises, $c^{\mathrm V\prime}(\theta)\le c^{\mathrm C\prime}(\theta)$. \\[2pt]
$W_i=W(\theta_i)\ge0$ & Longest time customer $i$ will queue at an FCS for a free gun (minutes). $W(\cdot)$ is non-increasing: a richer customer is less patient. $W_i=0$: customer $i$ only goes to a station with a free gun on arrival. \\[2pt]
$p^{\mathrm F},\ p^{\mathrm M}$ & Tariff the customer pays per kWh at an FCS and for the van service (Rs/kWh), $p^{\mathrm M}\ge p^{\mathrm F}$. \\
\bottomrule
\end{tabularx}}

\textbf{Shape of the sensitivity functions.} The model needs only the directions above. In the code each function is a straight line between its value at the lowest and at the highest income,
\[
\mu(\theta)=\mu^{0}+(\mu^{1}-\mu^{0})\,\theta,\qquad
c^{a}(\theta)=c^{a,0}+(c^{a,1}-c^{a,0})\,\theta,\qquad
W(\theta)=W^{0}+(W^{1}-W^{0})\,\theta,
\]
with $\mu^{0}\ge\mu^{1}>0$, $c^{a,1}\ge c^{a,0}\ge0$, $c^{\mathrm Q,0}>c^{\mathrm D,0}$, $c^{\mathrm Q,1}>c^{\mathrm D,1}$, $c^{\mathrm V,0}\le c^{\mathrm C,0}$, $c^{\mathrm V,1}-c^{\mathrm V,0}\le c^{\mathrm C,1}-c^{\mathrm C,0}$ and $W^{0}\ge W^{1}\ge0$. Because $\theta_i$ is data, $\mu_i$, $c^{a}_i$ and $W_i$ are numbers before the model is solved, so the model stays a linear MILP.

\textbf{Units.} Energy is in kWh, power in kW, distance in km and speed in km/h; every time quantity ($r_i$, $T$, $W_i$, $S_i$, $C_i$, $q_i$, $\tau$, $t$, big-$M$) is in minutes. Since kWh/kW and km/(km/h) give hours, charging and travel times are multiplied by 60. The values of time are per minute and $\mu_i$ is per rupee, so every term of $Z$ is in currency units. In the dataset the currency is Indian rupees (Rs).

\textbf{Values.} An ``$=$'' in the tables gives the value used in every file of the EV-Solomon dataset (\texttt{dataset\_ev\_solomon}); a value in parentheses after ``or'' applies to the 200-EV (Gehring--Homberger) files. $T$, $r_i$, $E_i$ and the locations come from the Solomon / Gehring--Homberger source files; all other values are our assumptions. The income parameters ($a_\theta$, $b_\theta$, the end points of $\mu$, $c^{a}$, $W$, and $p^{\mathrm F}$, $p^{\mathrm M}$) are kept symbolic here; the values used in the experiments are in the code (\texttt{code/f3\_params.json}).

\newpage
\subsection*{Decision Variables}
\begin{tabularx}{\textwidth}{@{}>{\raggedright\arraybackslash}p{4.6cm}X@{}}
\toprule
\textbf{Notation} & \textbf{Description} \\
\midrule
$y_{ig}\in\{0,1\}$ & 1 if EV $i$ is charged at gun $g\in G_i$; 0 otherwise. \\[2pt]
$z_{im}\in\{0,1\}$ & 1 if EV $i$ is served by van $m$; 0 otherwise. \\[2pt]
$x_{ijm}\in\{0,1\}$ & 1 if van $m$ traverses arc $(i,j)\in A$; 0 otherwise. \\[2pt]
$u_m\in\{0,1\}$ & 1 if van $m$ is dispatched; 0 otherwise. \\[2pt]
$o_{ij}\in\{0,1\}$ & For EV pairs $i<j$ sharing a gun, 1 if EV $i$ is charged before EV $j$; 0 otherwise. \\[2pt]
$S_i\ge0$ & Start time of EV $i$'s charging session, irrespective of the resource used (minute). \\[2pt]
$q^{\mathrm F}_i\ge0$ & Time EV $i$ queues at its FCS before charging starts (minutes); 0 for an MCS user. \\[2pt]
$q^{\mathrm M}_i\ge0$ & Time EV $i$ waits for its van before charging starts (minutes); 0 for an FCS user. \\
\bottomrule
\end{tabularx}

Short-hands:
\begin{gather*}
Y_{if}=\sum_{g\in G_f}y_{ig}\quad\text{(EV $i$ goes to FCS $f$)},\\
y_i=\sum_{g\in G_i}y_{ig},\quad z_i=\sum_{m\in M}z_{im}\quad\text{(FCS user / MCS user)},\\
C_i=S_i+\tau^{\mathrm F}_i\,y_i+\tau^{\mathrm M}_i\,z_i\quad\text{(completion time)}.
\end{gather*}
Because exactly one of $y_i$, $z_i$ equals 1, $C_i$ is the start time plus the duration on the chosen resource. The time from request to full charge splits into four parts:
\[
C_i-r_i\;=\;\underbrace{\sum_{f\in F_i}t_{if}Y_{if}}_{\text{driving}}\;+\;\underbrace{q^{\mathrm F}_i}_{\text{FCS queue}}\;+\;\underbrace{q^{\mathrm M}_i}_{\text{van wait}}\;+\;\underbrace{\tau^{\mathrm F}_i y_i+\tau^{\mathrm M}_i z_i}_{\text{charging}}.
\]
The money customer $i$ pays is
\[
\text{Pay}_i\;=\;\underbrace{c^{\mathrm{EV}}\sum_{f\in F_i}d_{if}Y_{if}+p^{\mathrm F}E_i\,y_i}_{\text{FCS user: driving + tariff}}\;+\;\underbrace{p^{\mathrm M}E_i\,z_i}_{\text{MCS user: tariff}}.
\]

%=====================================================================
\section{Objective}

\begin{equation}
\begin{aligned}
\min\ Z \;=\;& \underbrace{\sum_{i\in I}\mu_i\,c^{\mathrm{EV}}\sum_{f\in F_i}d_{if}Y_{if}}_{Z_1:\ \text{EV driving (perceived)}}
+\underbrace{c^{\mathrm M}\sum_{m\in M}\sum_{(i,j)\in A}d_{ij}x_{ijm}}_{Z_2:\ \text{van driving}}
+\underbrace{\sum_{m\in M}c^{\mathrm F}_mu_m}_{Z_3:\ \text{van dispatch}}\\
&+\underbrace{\pi^{\mathrm F}\sum_{i\in I}E_i\,y_i+\pi^{\mathrm M}\sum_{i\in I}\frac{E_i}{\eta}\,z_i}_{Z_4:\ \text{operator energy}}
+\underbrace{\sum_{i\in I}\mu_i\,E_i\big(p^{\mathrm F}y_i+p^{\mathrm M}z_i\big)}_{Z_5:\ \text{tariff (perceived)}}\\
&+\underbrace{\sum_{i\in I}\Big(c^{\mathrm D}_i\sum_{f\in F_i}t_{if}Y_{if}+c^{\mathrm Q}_i\,q^{\mathrm F}_i+c^{\mathrm C}_i\,\tau^{\mathrm F}_iy_i+c^{\mathrm V}_i\big(q^{\mathrm M}_i+\tau^{\mathrm M}_iz_i\big)\Big)}_{Z_6:\ \text{time (perceived)}}
\end{aligned}\label{eq:obj}
\end{equation}

$Z_2+Z_3+Z_4$ is the \textbf{operator's cost}, the same for every customer. $Z_1+Z_5+Z_6$ is the \textbf{customers' perceived cost}; every coefficient in it depends on the customer's income $\theta_i$.

$Z_1$ \textbf{EV driving.} An FCS user pays $c^{\mathrm{EV}}$ per km to drive to the station and feels it with weight $\mu_i$. A poorer customer minds a long drive more.

$Z_2$ \textbf{van driving.} Total km of all van routes (fuel/wear/driver time per km).

$Z_3$ \textbf{dispatch.} Fixed cost of every van that leaves the depot; with $Z_2$ this is the ``operating cost of MCS'' in the problem statement.

$Z_4$ \textbf{operator energy.} Grid price at an FCS versus van energy, which is dearer per kWh and loses $1-\eta$ in transfer. Without $Z_4$ the model would see van energy as free and overuse vans.

$Z_5$ \textbf{tariff.} What the customer pays per kWh, weighted by $\mu_i$. The van premium $(p^{\mathrm M}-p^{\mathrm F})E_i$ weighs heavily on a low-income customer and lightly on a high-income one.

$Z_6$ \textbf{time.} Each part of the time from request to full charge is priced at the customer's own value of time for that activity. An FCS user is with the car the whole time: driving ($c^{\mathrm D}_i$), queueing ($c^{\mathrm Q}_i>c^{\mathrm D}_i$) and charging ($c^{\mathrm C}_i$). An MCS user's car stays parked and the driver need not be present, so both the wait for the van and the charging are priced at $c^{\mathrm V}_i\le c^{\mathrm C}_i$. This is what the van sells to a time-sensitive customer: the charging time no longer costs the driver's time. If the nearest station A is full and station B has a free gun, the EV is sent to B as soon as $c^{\mathrm Q}_i\times$(queue at A) exceeds the extra driving cost and time to B; the richer the customer, the shorter the queue needed.

%---------------------------------------------------------------------
\subsection{How income enters the decision}

Fix the rest of the plan and compare the two ways of serving one EV $i$: at FCS $f$ with queue $q^{\mathrm F}_i$, or by a van with wait $q^{\mathrm M}_i$. From \eqref{eq:obj}, the extra cost of the van option is
\begin{align*}
\Delta_i(\theta_i)=\;&\underbrace{\mu(\theta_i)\big[(p^{\mathrm M}-p^{\mathrm F})E_i-c^{\mathrm{EV}}d_{if}\big]}_{\text{money}}
+\underbrace{\kappa_i}_{\text{operator}}\\
&-c^{\mathrm D}(\theta_i)\,t_{if}-c^{\mathrm Q}(\theta_i)\,q^{\mathrm F}_i-c^{\mathrm C}(\theta_i)\,\tau^{\mathrm F}_i+c^{\mathrm V}(\theta_i)\big(q^{\mathrm M}_i+\tau^{\mathrm M}_i\big),
\end{align*}
where $\kappa_i=(\pi^{\mathrm M}/\eta-\pi^{\mathrm F})E_i+c^{\mathrm M}\times$(extra van km) ($+\,c^{\mathrm F}_m$ if the van would not leave the depot otherwise) does not depend on income. The van is the better choice for EV $i$ when $\Delta_i<0$. How $\Delta_i$ changes with income:
\begin{align*}
\frac{d\Delta_i}{d\theta}=\;&\underbrace{\mu'(\theta)\big[(p^{\mathrm M}-p^{\mathrm F})E_i-c^{\mathrm{EV}}d_{if}\big]}_{\le 0\ \text{when the premium exceeds the drive cost}}\\
&\underbrace{-\,c^{\mathrm D\prime}(\theta)\,t_{if}-c^{\mathrm Q\prime}(\theta)\,q^{\mathrm F}_i-\big(c^{\mathrm C\prime}(\theta)\,\tau^{\mathrm F}_i-c^{\mathrm V\prime}(\theta)\,\tau^{\mathrm M}_i\big)}_{\le 0\ \text{since}\ c^{\mathrm V\prime}\le c^{\mathrm C\prime}\ \text{and}\ \tau^{\mathrm M}_i<\tau^{\mathrm F}_i}
\underbrace{+\,c^{\mathrm V\prime}(\theta)\,q^{\mathrm M}_i}_{\ge 0}.
\end{align*}
\begin{itemize}[itemsep=1pt]
  \item As long as the van wait $q^{\mathrm M}_i$ is short, $\Delta_i$ falls as income rises: the higher the income, the more attractive the van. Only a long van wait works the other way.
  \item The charging term drives this: at an FCS the driver spends $\tau^{\mathrm F}_i$ minutes at the rate $c^{\mathrm C}(\theta)$; with a van the car is charged in $\tau^{\mathrm M}_i<\tau^{\mathrm F}_i$ minutes at the lower driver-free rate $c^{\mathrm V}(\theta)$. The saving grows with income. If charging were priced the same for both resources, only the power difference $\tau^{\mathrm F}_i-\tau^{\mathrm M}_i$ would count and the van would lose most of its pull on high-income customers.
  \item With straight-line sensitivities $\Delta_i$ is linear in $\theta$, so for a given schedule there is one \emph{income threshold} $\theta^{*}_i$: below it EV $i$ goes to an FCS, above it a van serves it.
  \item The threshold falls (more customers get a van) when the van premium $p^{\mathrm M}-p^{\mathrm F}$ or the operator's extra cost $\kappa_i$ is small, when the EV is far from its station ($t_{if}$ large) and when FCS queues are long ($q^{\mathrm F}_i$ large, i.e.\ in the peak).
  \item Vans are scarce ($\sum_m Q_m$, battery, time). The model gives the van slots to the EVs with the most negative $\Delta_i$, which are mostly the high-income ones; the low-income customers absorb the FCS queues.
  \item The queue limit $W_i=W(\theta_i)$ adds a hard version of the same effect: in the peak a high-income customer cannot be left in a long queue at all, so they are the first to be moved to a free station or a van.
\end{itemize}
If $\mu$, every $c^{a}$ and $W$ are constant (no income effect), every customer is treated alike. The model is then formulation 2 without customer classes and without the preference level, plus the tariff terms, with van charging time priced at the driver-free rate.

%=====================================================================
\section{Constraints}

\subsection{Assignment}
\begin{equation}
y_i+z_i=\sum_{g\in G_i}y_{ig}+\sum_{m\in M}z_{im}=1\qquad\forall i\in I \label{eq:assign}
\end{equation}
\mdesc{Every request is served by exactly one resource: one gun of one reachable FCS, or one van. If EV $i$ reaches no FCS ($G_i=\emptyset$) it is forced onto a van, whatever its income.}

\subsection{Fixed charging stations (continuous time)}
\begin{equation}
S_i\ \ge\ r_i+\sum_{f\in F_i}t_{if}\,Y_{if}\qquad\forall i\in I \label{eq:fcsarr}
\end{equation}
\mdesc{An FCS user cannot start before it has been requested and has driven to its station. For an MCS user the sum is zero and this becomes $S_i\ge r_i$.}

\begin{align}
q^{\mathrm F}_i+q^{\mathrm M}_i &= S_i-r_i-\sum_{f\in F_i}t_{if}\,Y_{if} && \forall i\in I \label{eq:wait}\\
q^{\mathrm F}_i &\le W_i\,y_i && \forall i\in I \label{eq:waitF}\\
q^{\mathrm M}_i &\le T\,z_i && \forall i\in I \label{eq:waitM}
\end{align}
\mdesc{\eqref{eq:wait}: the time between arrival (FCS user) or request (MCS user) and the start of charging is the waiting time. \eqref{eq:waitF}: an FCS user queues at most $W_i$, its own income-dependent limit, so a long queue at one station forces the model to use another station or a van; an MCS user has no FCS queue. \eqref{eq:waitM}: only an MCS user can wait for a van; $T$ is a valid bound because $0\le S_i\le T$. With $q^{\mathrm F}_i,q^{\mathrm M}_i\ge0$, \eqref{eq:wait} implies \eqref{eq:fcsarr}, which is kept for readability.}

\begin{align}
S_j &\ge S_i+\tau^{\mathrm F}_i-M^{\mathrm F}_i\,(3-y_{ig}-y_{jg}-o_{ij}) && \forall g\in G,\ (i,j)\in P_g \label{eq:disj1}\\
S_i &\ge S_j+\tau^{\mathrm F}_j-M^{\mathrm F}_j\,(2-y_{ig}-y_{jg}+o_{ij}) && \forall g\in G,\ (i,j)\in P_g \label{eq:disj2}
\end{align}
\mdesc{No two EVs overlap on the same gun. If both $i$ and $j$ use gun $g$ ($y_{ig}=y_{jg}=1$), then either $o_{ij}=1$ and \eqref{eq:disj1} forces $j$ to start after $i$ has finished, or $o_{ij}=0$ and \eqref{eq:disj2} forces the reverse. If either is not on $g$ the big-$M$ switches both constraints off. With two guns per station this is exactly ``at most 2 EVs charging at an FCS at any moment'', without any time grid. $M^{\mathrm F}_i=T+\tau^{\mathrm F}_i$ is enough because $0\le S_i\le T$.}

\begin{equation}
\sum_{i\in I}\tau^{\mathrm F}_i\,y_{ig}\ \le\ T-\min_{i:\,g\in G_i}\Big(r_i+t_{if(g)}\Big)\qquad\forall g\in G^{+} \label{eq:gunload}
\end{equation}
\mdesc{Valid inequality (implied by the others, added to tighten the LP relaxation): the total charging time on one gun cannot exceed the time it is available. It is written only for guns in $G^{+}$: if no EV can reach gun $g$, the set $\{i: g\in G_i\}$ is empty, its minimum is undefined, and the gun has no $y_{ig}$ variables anyway, so nothing is lost.}

\begin{equation}
\sum_{i\in I}y_{ig_1}\ \ge\ \sum_{i\in I}y_{ig_2}\qquad\forall f\in F,\ G_f=\{g_1,g_2\} \label{eq:gunsym}
\end{equation}
\mdesc{Symmetry breaking: the two guns of a station are identical, so let gun 1 carry at least as many EVs as gun 2. Optional; it removes mirror solutions and speeds up branch-and-bound.}

\subsection{MCS routing}
\begin{align}
\sum_{j:(0,j)\in A}x_{0jm} &= u_m && \forall m\in M \label{eq:out}\\
\sum_{i:(i,n+1)\in A}x_{i,n+1,m} &= u_m && \forall m\in M \label{eq:in}\\
\sum_{i:(i,j)\in A}x_{ijm} &= z_{jm} && \forall j\in I,\ m\in M \label{eq:flowin}\\
\sum_{k:(j,k)\in A}x_{jkm} &= z_{jm} && \forall j\in I,\ m\in M \label{eq:flowout}
\end{align}
\mdesc{\eqref{eq:out}--\eqref{eq:in}: a dispatched van leaves the depot once and returns once; an unused van does neither. \eqref{eq:flowin}--\eqref{eq:flowout}: every EV assigned to van $m$ is entered once and left once by that van, i.e.\ visited exactly once; EVs not on $m$ are not touched by it. Together they make each van's arcs one path $0\to\cdots\to n+1$ (plus possible cycles, removed by \eqref{eq:seq}).}

\subsection{MCS capacity and battery}
\begin{align}
\sum_{i\in I}z_{im} &\le Q_m\,u_m && \forall m\in M \label{eq:cap}\\
\sum_{i\in I}\frac{E_i}{\eta}\,z_{im}+\beta\sum_{(i,j)\in A}d_{ij}\,x_{ijm} &\le B_m\,u_m && \forall m\in M \label{eq:batt}
\end{align}
\mdesc{\eqref{eq:cap}: a van serves at most $Q_m$ EVs, and none if it is not dispatched. \eqref{eq:batt}: the energy drawn from the van battery (energy delivered, grossed up for transfer losses, plus the energy to drive the whole route) must fit in its usable capacity.}

\subsection{MCS timing (one EV at a time, no subtours)}
\begin{align}
S_j &\ge t_{0j}\,x_{0jm} && \forall j\in I,\ m\in M \label{eq:depart}\\
S_j &\ge S_i+\tau^{\mathrm M}_i+t_{ij}-M^{\mathrm M}_{ij}\,(1-x_{ijm}) && \forall i,j\in I,\ i\ne j,\ m\in M \label{eq:seq}\\
S_i+\tau^{\mathrm M}_i+t_{i,n+1} &\le T+(\tau^{\mathrm M}_i+t_{i,n+1})\,(1-x_{i,n+1,m}) && \forall i\in I,\ m\in M \label{eq:return}
\end{align}
\mdesc{\eqref{eq:depart}: the first EV on a route cannot be started before the van (leaving at 0) has driven to it. \eqref{eq:seq}: if van $m$ goes from $i$ to $j$, it must finish charging $i$ and drive to $j$ before starting $j$. This both prevents a van from charging two EVs at once and, because $\tau^{\mathrm M}_i+t_{ij}>0$, makes start times strictly increase along a route so no cycle can close (MTZ-type subtour elimination). \eqref{eq:return}: the van is back at the depot by $T$.}

\subsection{Time windows}
\begin{align}
S_i &\ge r_i && \forall i\in I \label{eq:release}\\
C_i &\le T && \forall i\in I \label{eq:horizon}
\end{align}
\mdesc{\eqref{eq:release}: no charging before the request (a van that arrives early waits). It is implied by \eqref{eq:fcsarr} and kept for readability. \eqref{eq:horizon}: every EV is fully charged within the peak, which guarantees ``all requests satisfied''.}

\subsection{Van symmetry and domains}
\begin{align}
u_m &\ge u_{m+1} && m=1,\dots,|M|-1 \label{eq:vansym}\\
y_{ig},\ z_{im},\ x_{ijm},\ u_m,\ o_{ij}\in\{0,1\}, & \qquad S_i,\ q^{\mathrm F}_i,\ q^{\mathrm M}_i\ge0 \label{eq:dom}
\end{align}
\mdesc{\eqref{eq:vansym}: vans are identical, so use them in index order (van 2 only if van 1 is used, etc.). \eqref{eq:dom}: variable domains.}

%=====================================================================
\section{Changes with respect to formulation 2}
\begin{enumerate}[itemsep=2pt]
  \item \textbf{No a-priori classes.} The classes $K=\{\mathrm L,\mathrm{Mid},\mathrm H\}$, the sets $I^{k}$, the class map $k(i)$, the switch count $V$ and the lexicographic level are removed. The model has one objective and is solved once.
  \item \textbf{Continuous income attribute} $\theta_i\in[0,1]$ per EV, drawn from a distribution, with the sensitivity functions $\mu(\theta)$ (money), $c^{a}(\theta)$ (time, four activities) and $W(\theta)$ (queue tolerance). The class rates $c^{a}_k$ of formulation 2 become $c^{a}_i=c^{a}(\theta_i)$.
  \item \textbf{Perceived money cost.} Customer tariffs $p^{\mathrm F}\le p^{\mathrm M}$ and the money sensitivity enter through $Z_1$ (EV driving, now weighted by $\mu_i$) and the new term $Z_5$ (tariff). The time term (formulation 2's $Z_5$) is now $Z_6$.
  \item \textbf{Queue limit per customer.} $W_{\max}$ becomes $W_i=W(\theta_i)$; constraints (7) and (9) of formulation 2 merge into \eqref{eq:waitF}, $q^{\mathrm F}_i\le W_i\,y_i$.
  \item \textbf{Charging time priced by where the driver is.} Charging at an FCS costs $c^{\mathrm C}_i$ per minute (driver present); charging by a van costs $c^{\mathrm V}_i$ per minute, like the wait for the van (car parked, driver free). In formulation 2 both charging times used $c^{\mathrm C}_k$.
\end{enumerate}

\end{document}
